\documentclass[tikz,border=10pt]{standalone}
\begin{document}
\begin{tikzpicture}[scale=1.0]
    % Define coordinates
    \coordinate (O) at (0,0);
    \coordinate (B) at (-4,0);
    \coordinate (E) at (4,0);
    % A and F at 120 and 60 degrees
    \coordinate (A) at (-2, 3.464);
    \coordinate (F) at (2, 3.464);
    % C and D at 210 and 330 degrees
    \coordinate (C) at (-3.464, -2);
    \coordinate (D) at (3.464, -2);
    % H is intersection of BA and EF extended. 
    % Line BA: y = sqrt(3)(x+4). Line EF: y = -sqrt(3)(x-4).
    % Intersection at x=0, y=4*sqrt(3) approx 6.928
    \coordinate (H) at (0, 6.928);
    % G is intersection of CD extended and vertical line from E
    % Line CD is y = -2. Vertical line at E is x = 4.
    \coordinate (G) at (4, -2);

    % Draw Circle
    \draw (O) circle (4);

    % Draw Lines
    \draw (H) -- (B); % Passes through A
    \draw (H) -- (E); % Passes through F
    \draw (A) -- (F);
    \draw (B) -- (F);
    \draw (B) -- (D);
    \draw (B) -- (C);
    \draw (B) -- (E); % Diameter
    \draw (C) -- (G); % Passes through D
    \draw (D) -- (E);
    \draw (E) -- (G);

    % Right angle mark at G
    \draw (G) ++(0, 0.3) -- ++(-0.3, 0) -- ++(0, -0.3);

    % Draw Points (small circles if needed, but labels are enough usually. 
    % The image has a dot for O. Others are intersections.)
    \fill (O) circle (2pt);

    % Labels
    \node[above] at (H) {$H$};
    \node[above left] at (A) {$A$};
    \node[above right] at (F) {$F$};
    \node[left] at (B) {$B$};
    \node[right] at (E) {$E$};
    \node[below] at (O) {$O$};
    \node[below left] at (C) {$C$};
    \node[below] at (D) {$D$};
    \node[below right] at (G) {$G$};

\end{tikzpicture}
\end{document}
